\par\noindent\textbf{(1)} 相干振幅点扩散函数为瞳函数 \(P(\boldsymbol r)\) 的Fourier变换，非相干强度点扩散函数 \(h=|\mathcal F P|^2\)。光学传递函数 \(\mathrm{OTF}(\boldsymbol\nu)=\mathcal F h\) 因Wiener–Khinchin定理等于瞳函数的归一自相关。
\par\noindent\textbf{(2)} 两个半径 \(R\) 圆瞳平移距离 \(s\) 的交叠面积为 \(A_\cap=2R^2\arccos(s/(2R))-(s/2)\sqrt{4R^2-s^2}\)。用原面积 \(\pi R^2\) 归一，得到 \(\mathrm{MTF}(\rho)=\frac2\pi[\arccos\rho-\rho\sqrt{1-\rho^2}]\)，其中 \(\rho=s/(2R)\in[0,1]\)。
\par\noindent\textbf{(3)} 对于像面坐标 \(x=f\theta\)，瞳孔半径 \(R=D/2\) 对应相干振幅空间频率支撑半径 \(\nu_{\mathrm{coh}}=R/(\lambda f)=D/(2\lambda f)\)，非相干OTF的自相关支撑直径为 \(\nu_{\max}=2R/(\lambda f)=D/(\lambda f)\)。代入 \(\rho=\nu/\nu_{\max}\)，当 \(\nu\ge\nu_{\max}\) 时两个圆瞳无交集，\(\mathrm{MTF}=0\)。不能误把直径与半径所定义的截止频率再多乘一倍。
\par\noindent\textbf{(4)} 在 \(\rho=0\) 时 \(\mathrm{MTF}=1\)，对低频展开为 \(1-(4/\pi)\rho+O(\rho^3)\)。像差引入瞳函数的相位变化，虽然支撑截止频率不变，但自相关的模可能显著下降；不能只看几何孔径而断言实际系统达到衍射极限对比度。
